\chapter{דרישות המערכת}\label{ch:req}

הדרישות מנוסחות לפי \en{ISO/IEC/IEEE 29148}~\cite{iso29148}: כל דרישה חד משמעית, ניתנת לבדיקה, ומקבלת מזהה קבוע. העדיפות נקבעת בשיטת \en{MoSCoW}: \textbf{M}, חובה, \textbf{S}, רצוי מאוד, \textbf{C}, אם יתאפשר. שיטת האימות: \textbf{T}, בדיקה, \textbf{A}, ניתוח, \textbf{D}, הדגמה, \textbf{I}, בחינה. כל דרישה בעדיפות \textbf{M} מקושרת לבדיקה בנספח~\ref{app:trace}.

\section{דרישות פונקציונליות}
\small

ראו טבלה~\ref{tab:freq}.

\begin{longtable}{@{}l p{9.3cm} c c@{}}
\caption{דרישות פונקציונליות}\label{tab:freq}\\
\toprule
\textbf{מזהה} & \textbf{דרישה} & \textbf{עדיפות} & \textbf{אימות} \\
\midrule
\endfirsthead
\toprule
\textbf{מזהה} & \textbf{דרישה} & \textbf{עדיפות} & \textbf{אימות} \\
\midrule
\endhead
\bottomrule
\endfoot
\multicolumn{4}{@{}l}{\textit{קליטה}}\\
\req{FR-01} & המערכת תקלוט תמונה בהעלאה ידנית, בגרירה למערכת ניהול התוכן, ובממשק \en{API}. & M & T \\
\req{FR-02} & המערכת תתמוך בפורמטים \en{JPEG}, \en{PNG}, \en{WebP}, \en{HEIC} עד \num{50} מגה בייט לקובץ. & M & T \\
\req{FR-03} & המערכת תשמור את קובץ המקור ללא שינוי ותחשב לו טביעת \en{SHA-256}. & M & I \\
\multicolumn{4}{@{}l}{\textit{שכבת מקור}}\\
\req{FR-04} & המערכת תקרא ותאמת תעודת \en{C2PA} ותציג את שרשרת העריכות. & M & T \\
\req{FR-05} & המערכת תזהה תעודה פגומה, שפג תוקפה או שנחתמה בתעודה שבוטלה. & M & T \\
\req{FR-06} & המערכת תחפש סימני מים מוכרים של מחוללים (לפי רשימה מתעדכנת). & S & T \\
\multicolumn{4}{@{}l}{\textit{שכבת זיהוי}}\\
\req{FR-07} & המערכת תריץ לפחות ארבעה גלאים בלתי תלויים ותשקלל אותם לציון יחיד \num{0} עד \num{100}. & M & T \\
\req{FR-08} & הציון יכויל כך שציון \num{80} ומעלה יתאים לשיעור זיוף בפועל של \num{80\%} לפחות במאגר הכיול. & M & A \\
\req{FR-09} & המערכת תציג מפת חום שמסמנת את האזורים שתרמו לציון. & S & D \\
\multicolumn{4}{@{}l}{\textit{שכבת הקשר}}\\
\req{FR-10} & המערכת תבצע חיפוש תמונה הפוך ותציג את ההופעה המוקדמת ביותר שנמצאה ברשת. & M & T \\
\req{FR-11} & המערכת תשווה מטא דאטה (תאריך, מיקום, מכשיר) לטענה בכיתוב התמונה ותסמן סתירות. & S & T \\
\req{FR-12} & המערכת תציע לבודק רשימת בדיקות מודיעין גלוי לפי סוג האירוע (אסון, הפגנה, מלחמה). & C & D \\
\multicolumn{4}{@{}l}{\textit{החלטה ותיוג}}\\
\req{FR-13} & המערכת תסווג כל תמונה לאחת משלוש החלטות: ״מאומת״, ״לא ודאי״, ״חשד גבוה״. & M & T \\
\req{FR-14} & תמונה בסיווג ״לא ודאי״ או ״חשד גבוה״ לא תפורסם כתיעוד בלי אישור עורך בכיר. & M & T \\
\req{FR-15} & המערכת תצמיד לכל תמונה שפורסמה תווית לקורא לפי \en{TRUTH-IMG}, עם קישור לתעודת המקור. & M & I \\
\req{FR-16} & אחרי אישור, המערכת תחתום על התמונה בתעודת \en{C2PA} של מערכת החדשות. & S & T \\
\multicolumn{4}{@{}l}{\textit{תיעוד ודיווח}}\\
\req{FR-17} & כל החלטה תירשם ביומן בלתי ניתן לשינוי: מי, מתי, ציון, החלטה ונימוק. & M & I \\
\req{FR-18} & המערכת תפיק דוח שקיפות חודשי למערכת החדשות ולציבור. & M & D \\
\req{FR-19} & לוח הבקרה יציג: היסטוגרמת ציונים, עקומת \en{ROC} חודשית, וגרף עמודות של החלטות לפי מחלקה. & M & D \\
\req{FR-20} & ממשק \en{API} ״ציון שלמות״ ימסור לרשתות פרסום ציון מצרפי לאתר, בלי פרטי תמונות. & S & T \\
\req{FR-21} & המערכת תסרוק מחדש את ארכיון \num{90} הימים האחרונים כשמתווסף גלאי חדש. & C & T \\
\end{longtable}
\normalsize

\section{דרישות ביצועים}

ראו טבלה~\ref{tab:nfr}.

\begin{table}[!htbp]\centering
\caption{דרישות ביצועים, אבטחה ואילוצים}\label{tab:nfr}
\small
\begin{tabularx}{\linewidth}{@{}l Y c c@{}}
\toprule
\textbf{מזהה} & \textbf{דרישה} & \textbf{עדיפות} & \textbf{אימות} \\
\midrule
\req{PR-01} & זמן תגובה אוטומטי: חציון $\leq$ \num{60} שנ׳, אחוזון \num{95} $\leq$ \num{180} שנ׳. & M & T \\
\req{PR-02} & קיבולת: \num{2} מיליון תמונות בחודש, מוכחת במבחן עומס (הרבה מעל נפח הפיילוט), עם יכולת הרחבה לפי \num{25} בלי שינוי ארכיטקטורה. & M & A \\
\req{PR-03} & זמינות \num{99.9\%} בחודש (עד \num{44} דקות השבתה). & M & A \\
\req{PR-04} & דיוק: שיעור גילוי של \num{90\%} לפחות ושיעור התרעות שווא של \num{2\%} לכל היותר, על מאגר \en{VL-Wild}. & M & T \\
\req{PR-05} & עמידות: ירידה של עד \num{5} נקודות בשיעור הגילוי אחרי דחיסת \en{JPEG} באיכות \num{70} והקטנה ל-\num{50\%}. & M & T \\
\midrule
\req{SR-01} & הצפנה במעבר (\en{TLS 1.3}) ובמנוחה (\en{AES-256}). & M & I \\
\req{SR-02} & הרשאות לפי תפקיד: צלם, עורך, עורך בכיר, מנהל, מבקר. & M & T \\
\req{SR-03} & תמונות שלא פורסמו יימחקו אחרי \num{30} יום, אלא אם נקבע אחרת בחוזה. & M & I \\
\req{SR-04} & ציות ל-\en{GDPR} ולחוק הגנת הפרטיות הישראלי; אחסון הנתונים באזור שבחר הלקוח. & M & I \\
\midrule
\req{CR-01} & ממשק משתמש בעברית, אנגלית, ערבית וצרפתית, כולל תמיכה מלאה בכתיבה מימין לשמאל. & M & D \\
\req{CR-02} & תוספים מוכנים ל-\en{WordPress} ול-\en{Arc XP}; \en{API} מתועד (\en{OpenAPI 3.1}) לכל מערכת ניהול תוכן אחרת. & M & D \\
\req{CR-03} & עמידה בהנחיות סעיף \num{50} לחוק הבינה המלאכותית האירופי לגבי גילוי ״דיפ פייק״. & M & A \\
\req{CR-04} & נגישות ממשק לפי \en{WCAG 2.2} ברמה \en{AA}. & S & I \\
\bottomrule
\end{tabularx}
\end{table}

\section{דרישות ממשק}
המערכת מחוברת לחמישה גורמים חיצוניים (איור~\ref{fig:context} בפרק~\ref{ch:design}): מערכת ניהול התוכן של הלקוח, אפליקציית הצלם, שירותי חיפוש תמונה הפוך, רשימות ביטול של תעודות \en{C2PA}, ורשתות פרסום. כל ממשק מוגדר בחוזה \en{API} עם גרסה, ושינוי שובר תאימות מחייב הודעה של \num{90} יום מראש.

\section{דרישות התקן האתי}
לצד הדרישות הטכניות, הפרויקט מחויב לתוצר נורמטיבי: מסמך \en{TRUTH-IMG} בגרסה \num{1.0}, שכולל עשרה כללים (פרק~\ref{ch:design}, סעיף~\ref{sec:truthimg}), נוהל ביקורת שנתי, וטופס הסמכה. התקן ייבחן בסבב הערות ציבורי של \num{30} יום מול לפחות שלוש מערכות חדשות וגוף אתיקה אחד לפני פרסומו.
